长期记忆的困难在于：历史不断增长，既要保持易于检索，又不能丢失后续作答所需的精确措辞。把历史装入有限的上下文窗口，本质上是一个内存管理问题：进入窗口的内容应当被寻址，而不是整块搬运。Deep-Dream 提供的正是寻址结构：原始文档与精确段落构成事实记录；概念、关系、观测与版本链接负责指明去哪里找；提问时，智能体只打开并核对回答该问题所需的源文。若早期提及无法确定所指实体，系统保留多个候选，后续观测可以重定向链接而不改写原文。我们在 MemoryAgentBench、LongMemEval、LoCoMo 和 MEME 上评测该设计。在 MemoryAgentBench 上，同一模型、同一评分器下，单次检索结果作答得 0.333，只读记忆工具作答得 0.400，智能体亲自打开的源文作答得 0.710，增益集中在多次提及与更新类问题；同模型全量语料参考界定了选择性阅读尚未覆盖的部分。Mem0 与 Zep 在 LoCoMo 和 LongMemEval 上的已发表分数使用不同的评审模型与评分规则版本，仅作背景参照，不是受控比较。在配对运行中，先记录观测、待更多上下文可用后再合并相容身份的修订写入实现，使每份文档的写入请求减少 46--65\%，而抽样质量相近。
